\chapter{Marketing}
\section{Analyse Pestel}
L'analyse marketing a été structurée en deux grandes parties. D'une part, nouvs avons collecté un maximum de données concernant les besoins et attentes des patient en commencant par une analyse Pestel détaillée. Cette analyse nous a permis d'établir le contexte de marché. Les grands éléments à retenir pour cette analyse sontes les suivantes
 
\begin{itemize}
    \item Politique: Le développement du liner dépendra des subsides accordées par la région Wallonne. Ces subsides seront accordés suite à une compétition de financement publique auquel différents instituts (universités/hautes écoles) présenteront leurs projets. Le dossier consituté pour la région devra donc être particulièrement sourcé et solide
    \item Economique: Le marché mondial des prothèses de membres inférieurs est en forte croissance et devrait atteindre 2,6 milliards USD en 2030 contre 1,7 milliards USD en 2024. En plus d'un marché croissant, le système de remboursement INAMI offre un cadre favorable pour un tel marché avec des remboursements allant jusqu'à 2696,83€ pour un liner adaptatif. 
    \item Socioculturel: Les patients amputés font fassent à des douleurs et gènes au cours de la journée. La dilatation et retraction du moignon est incontrolable et oblige les patients à faire des ajustements pour limiter les gènes. Un dispositif comme un liner adaptatif semble être une solution novatrice jamais expérimentée auparavant.
    \item Technologique: Le futur liner utilisera des technologies de pointeavec capteurs de pression piézorésistifs ou capacitifs intégrés avec électrodes textiles soiples et lavables. Ces solutions sont développés par des entreprises et instituts européenne comme AG-tech (France) et les universités de Nantes et de Louvain (Louvain Bionics). Les défos technologique sont cependant majeurs avec une gestion de batterie qui devra être cruciale et une miniaturisation des composants qui ne doivent créer aucune gène chez le patient.
    \item Environnemental: Les ressources utilisées devront être le plus local possible car la région Wallonne y fait particulièrement attention. Malgré cela, certains composants ne pourront pas être exclusivement produits en Belgique. De plus, les contraintes de stockages des produits biomédicaux doivent respecter des températures et taux d'humidités précis.
    \item Légal: Le liner nouvellement développée devra respecter des normes et caractéristiques très précises afin d'être classifié comme un produit personnalisé. La position des capteurs de pression doit être à tout prix personnalisée pour que le produit puisse être vendu pour ~2700€. Les données collectées chez les patients doivent évidemment rester confidentielles en vertue de la protection des données personnelles (RGPD).
\end{itemize}
 
Cette analyse Pestel recensent le cadre dans lequel évoluera l'entreprise nouvellement crée et sera utilisée par la suite pour identifier les opportunités et menaces de la matrice SWOT.
\section{Exploration du marché (FIELD)}
 
L'exploration du marché a été réalisée en deux grandes parties, la première étant la collecte de donnée prilaire qualitatives. Pour réaliser cette collecte, un sondage de 33 questions a été réalisé et partager avec Mr.Fischer qui a formulé différente recommandations pour préciser les questions et assurer que le sondage collecte toutes les données nécéssaires auprès des patients. Une fois le sondage finalisé, celui-ci a été transmis une nouvelle fois et est actuellemnt diffusé auprès des clients de Tecfit qui donne leurs avis, expriment leurs espoirs/craintes par rapport à un liner automatique et soumettent leurs informations personnelles pour une interview en présentiel par la suite.  
 
\section{Exploration du marché (DESK)}
Pour les données secondaires, un grand nombre de sources ont été explorées afin de déterminer la taille du marché des liners en Belgique et en Europe de l'ouest. La première tâche a été de connaître l'évolution des amputation majeurs des membres inférieurs en Belgique. Les amputations majeurs sont celles réalisées au dessus de la cheville (tibia, genou, cuisse) et constitue le marché des patients susceptibles d'utiliser un liner. Deux sources ont été consultées, et fournissent l'évolution des amputations entre 2009 à 2013 et 2009 à 2018 pour le deuxième article. Une droite de régression a ensuite été tracée pour estimer les amputations de cette année et voici les résultats obtenus:
 
\begin{figure}[H]
    \centering
    \includegraphics[width=1\linewidth]{Fig_Marketing/Projection_2009_2018.png}
    \caption{Projection du nombre d'amputations}    \label{fig:Projetction_2009_2018}
\end{figure}


Nous constatons un taux d'amputations qui est en baisse à raison de 25 amputations en moins par an. Naturellement, ces chiffres se base sur la décénie passée et un plateau sera atteint car il n'est pas possible de ne plus avoir recourt à l'amputation, et ce malgré l'avancée de la médecine. Nous estimons donc que le nombre d'amputations pour 2025 se situe entre 950 et 1000. Ceci nous informe sur la tendance du marché, mais pas sur le nombre total de patients susceptiblent d'avoir recourt à un liner.
Il est très dur de connaître la taille du marché au patient près, un chiffre datant de 2022 nous informe qu'il y a +-7700 amputées en Belgiques (toutes amputations confondues). Ce chiffre se base sur une extrapolation du nombre d'amputés au Pays-Bas recalculé pour la population belge.
Le Pays-Bas étant un pays limitrophe de la Belgique avec un indice de développement humain similaire, ce chiffre de 7700 constitue une estimation solide. Sur base des différents articles collectés (Allemagne, Belgique, France, Pays-bas,...), nous constatons que les amputation des membre inférieurs représentent une minorité de cas et compte pour 5\% du nombre d'amputations total.
7300 est donc le nombre d'amputés de la jambe en Belgique. Sur base de ces même documents, le nombre d'amputations majeurs de la jambe par rapport au nombre d'amputations mineurs varie de 28\% jusqu'à 40 \% en fonction des sondages et des années. 33\% est donc un pourcentage moyen. Le nombre de patient avec amputations majeurs de la jambe en Belgique se situe entre 2400 et 2500 selon nos calculs (tous types confondus: diabétiques type 1, diabétiques type 2, accidents de travail).
Cela représente donc le marché maximal pouvant être atteint en Belgique. \\
 
Ce chiffre de 7700 bien que cité dans beaucoup de sources n'est pas justifié et aucun article ne semble expliquer comment il a été obtenu. Nous avons effectué une deuxième estimation du marché en se basant sur le taux de décès par an dans la population touchée par une amputation majeure. Plusieurs sources indiqiquent un taux de décès à 5 ans de 69\% pour les personnes diabétiques et 60\% pour les personnes non diébatiques. Cela nous donne un taux de survie à 5 ans de 35\%. 
\\
Taux de survie à 5 ans:
\begin{equation}
S_5 = 1 - 0,65 = 0,35
\end{equation}
Taux de mortalité par an:
\begin{equation}
S_5 = (1 - M)^5 => M = 1 - 0,35^\frac{1}{5} = 0,19
\end{equation}
 
Rapporté sur une année, le taux de décès est de 19\% de la population avec amputation majeure. Avec les nouveaux amputées précédemment calculer avec la régression précédente, voici les chiffres projetés pour 2025:
 
\begin{figure}[H]
    \centering
    \includegraphics[width=1\linewidth]{Fig_Marketing/Chiffres_Decenies.png}
    \caption{Projection du nombre d'amputés sur la décénie}    \label{fig:Chiffres_Decenies}
\end{figure}
 
L'ordre de grandeur par rapport au calcul prédédent est similaire, mais la valeur du marché total passe du simple au double. Pour conserver une approche conservatrice dans l'estimation du marché, la première valeur a été retenue. Nous etimerons donc que le nombre de personnes avec une amputation majeures en Belgique est de 2500.
 
\section{Besoin en flux de trésorerie}
Suite à une estimation des coûts fixes et variables sur 5 ans effectuée dans la partie logistique, une première courbe des besoins en flux de trésorerie a été réalisée. Cette courbe a pour but de mettre en évidence le besoin maximal en trésorerie et quand ce besoin maximal est atteint. En considérant une évolution linéaire (les patients changent de liner tous les ans) voici les résultats obtenus:
 
\begin{figure}[H]
    \centering
    \includegraphics[width=0.8\linewidth]{Fig_Marketing/Courbe_flux_tresorerie.png}
    \caption{Projection du nombre d'amputés sur la décénie}    \label{fig:Courbe_flux_tresorerie}
\end{figure}
 
\begin{itemize}
    \item Le besoin maximal est de 600000€ a lieu au bout de 25 mois 
    \item La rentabilité est atteinte au bout de 50 semaines
    \item Le nombre de pacients est croissant à raison de 3 nouveaux liners fabriqués tous les deux mois
\end{itemize}


